% =====================================================================
\BookChapter[C++]{C++: Abstraction Without Losing Speed}
\label{ch:cpp}

\begin{researchquestion}
How can one language be used for both a microcontroller and a million-line simulation package?
\end{researchquestion}
\OrientationNote
% =====================================================================

\begin{factsheet}{C++}
First appeared & 1985 (``C with Classes'' from 1979); standards C++98, 11, 14, 17, 20, 23\\
Designed by & Bjarne Stroustrup, Bell Labs\\
Paradigms & procedural, object-oriented, generic (templates), functional elements\\
Typing & static, mostly strong\\
Execution & compiled to native code\\
File extensions & \file{.cpp}, \file{.cc}, \file{.cxx} (source); \file{.hpp}, \file{.h} (headers)\\
Compilers & GCC (\code{g++}), Clang (\code{clang++}), Microsoft Visual C++, Intel\\
Used for & simulation software (FEM, CFD, MD), games and graphics, CAD, robotics,
  trading systems, browsers, machine-learning frameworks, embedded systems\\
Famous software & OpenFOAM, LAMMPS, deal.II, Gmsh, ROS, Unreal Engine, Chrome, the cores
  of PyTorch and TensorFlow, Eigen, Blender, much of Ansys, Abaqus and COMSOL\\
\end{factsheet}

C++ takes C and adds abstractions: classes, templates, a rich standard library, and
automatic resource management, while keeping C's speed and its access to the hardware. It
is the language of large, performance-critical engineering software. When you run a
finite-element analysis, an open-source CFD case, a robot or a modern video game, there is
very likely C++ underneath.

\LearningOutcomes{Replace error-prone C constructs with their safer C++ equivalents; pass objects by reference; write a small class hierarchy with virtual functions; explain how resources are released automatically; use standard containers and algorithms; and name the libraries that engineers use for numerical work in C++.}

\section{History and philosophy}

\begin{history}
Bjarne Stroustrup, a Danish computer scientist, had used the language Simula during his
PhD at Cambridge to simulate distributed computer systems. Simula's classes made the
program easy to organise, but it was far too slow. In 1979, at Bell Labs, he began adding
Simula-style classes to C to get both organisation and speed. ``C with Classes'' was
renamed C++ in 1983 (\code{++} is C's increment operator: C++ is ``one more than C''). The
first international standard appeared in 1998. Since C++11 the language has been updated
every three years, and modern C++ is a much safer and more pleasant language than the C++
of the 1990s.
\end{history}

Stroustrup's guiding principle is the \term{zero-overhead principle}: you do not pay
(in speed or memory) for features you do not use, and what you do use is as efficient as
if you had written it by hand. This makes C++ suitable both for high-level design and for
the innermost loops of a solver. The price is complexity: C++ is a very large language, and
nobody uses all of it. A beginner should learn a modern, safe subset.

\section{Getting started}

Any C compiler distribution also includes a C++ compiler: \code{g++} or \code{clang++} on
Linux and macOS, Visual Studio on Windows. VS Code, CLion and Qt Creator are popular IDEs.
\codefile{C++}{cpp/hello.cpp}{hello world}
\begin{lstlisting}[style=shell]
$ g++ -std=c++20 -Wall -Wextra -o hello hello.cpp
$ ./hello
\end{lstlisting}
\outfile{cpp/hello.out}
Instead of \code{printf}, C++ uses \term{streams}: \code{std::cout} (``character
output'') with the insertion operator \code{<<}. The prefix \code{std::} says that
\code{cout} belongs to the \term{namespace} of the standard library; namespaces prevent
name clashes between libraries, much like \py's modules. In \code{main}, the final
\code{return 0;} may be omitted.

Almost every valid C program is also a valid C++ program, so everything in
Chapter~\ref{ch:c} still applies. But good C++ code looks quite different from C code.

\section{From C to modern C++}

\subsection*{Use the safer building blocks of the standard library}

Modern C++ replaces the most error-prone C features with safer library types.
\begin{center}
\small
\begin{tabularx}{\linewidth}{@{}l>{\raggedright\arraybackslash}X>{\raggedright\arraybackslash}X@{}}
\toprule
\textbf{Task} & \textbf{C style (avoid in C++)} & \textbf{Modern C++}\\
\midrule
dynamic array & \code{malloc}/\code{free}, raw pointer & \code{std::vector<double>}\\
fixed array & \code{double a[3]} & \code{std::array<double, 3>}\\
text & \code{char name[32]} & \code{std::string}\\
own an object & \code{new}/\code{delete} & \code{std::unique\_ptr}, \code{std::make\_unique}\\
constants & \code{\#define N 40} & \code{constexpr int N = 40;}\\
type conversion & \code{(int) x} & \code{static\_cast<int>(x)}\\
output & \code{printf} & \code{std::cout}, or \code{std::print} (C++23)\\
\bottomrule
\end{tabularx}
\end{center}

\code{std::vector} is the workhorse. It knows its own size, grows when needed, and releases
its memory automatically when it goes out of scope:
\begin{lstlisting}[language=C++]
#include <vector>
std::vector<double> T = {293.1, 305.6, 299.4};
T.push_back(301.2);             // like Python's append
T.size();                       // 4
T[0];                           // no bounds check (fast)
T.at(10);                       // bounds check: throws std::out_of_range
for (double x : T) { /* ... */ } // range-based for loop, like Python's for
\end{lstlisting}

\subsection*{Pass large objects by reference, not by copy}

C++ adds \term{references}, a safer alternative to pointers for passing arguments. A
parameter declared as \code{double\&} is another name for the caller's variable; a
parameter declared as \code{const std::vector<double>\&} gives the function read-only access
to the caller's vector without copying it:
\begin{lstlisting}[language=C++]
void scale(std::vector<double>& x, double factor)   // modifies the caller's vector
{
    for (double& xi : x) xi *= factor;
}

double norm(const std::vector<double>& x)           // reads it without copying
{
    double s = 0.0;
    for (double xi : x) s += xi * xi;
    return std::sqrt(s);
}
\end{lstlisting}
Passing large objects by \code{const} reference is one of the first habits of a C++
programmer.

\subsection*{Let the compiler deduce types, and write small functions in place}

\code{auto} lets the compiler work out the type of a variable from its initial value.
\term[lambda expressions]{Lambda expressions} define small unnamed functions inline, like \py's \code{lambda}
but with any number of statements. The brackets \code{[\,]} say which surrounding variables
the lambda may use (\emph{capture}): \code{[=]} copies them, \code{[\&]} refers to them.
\begin{lstlisting}[language=C++]
auto k = 0.15;                                    // double
auto decay = [k](double t) { return std::exp(-k * t); };
std::cout << decay(2.0) << '\n';                  // 0.740818
\end{lstlisting}

\section{Case study: the falling skydiver}

\codefile{C++}{cpp/drag.cpp}{Euler solution of the drag equation}
\outfile{cpp/drag.out}
Compared with the C version, there are no global variables and no fixed-size arrays: the
parameters are \code{const} local variables, the right-hand side is a lambda that captures
them, and the results go into \code{std::vector}s whose size is computed at run time.
We kept \code{std::printf} because its format strings are compact; C++23's \code{std::print}
offers the same convenience with \py-like format strings, \code{std::print("\{:6.1f\}",
t)}.

\section{Classes and object-oriented programming}

A \code{class} in C++ bundles data members and member functions. Members are
\code{private} (accessible only inside the class) or \code{public}. A \term{constructor}
initialises a new object. Inheritance lets a class extend another one, and \code{virtual}
functions let the program choose the right version of a function at run time
(\term{polymorphism}).

The following example models beam cross-sections. The bending stiffness of a beam is $EI$,
where $I$ is the second moment of area of its cross-section. A base class \code{Section}
declares what every section must provide; three derived classes implement it.
\codefile{C++}{cpp/sections.cpp}{classes, inheritance and polymorphism}
\outfile{cpp/sections.out}
Points to notice:
\begin{itemize}
\item \code{Section} is an \term{abstract class}: its functions are declared \code{= 0}, so
  it cannot be instantiated. It defines an \emph{interface}.
\item \code{override} asks the compiler to check that a function really overrides a
  virtual function of the base class --- it catches spelling mistakes.
\item \code{const} after a member function promises that it does not modify the object.
\item The loop in \code{main} calls \code{s->area()} without knowing what kind of section
  \code{s} is. The tube, which has a much smaller area than the solid circle but a similar
  second moment, has the largest $I/A$ ratio: this is why hollow sections are efficient for
  bending and torsion.
\item \code{std::unique\_ptr} owns each object and deletes it automatically. There is no
  \code{delete} anywhere in the program.
\end{itemize}

\begin{keyidea}[Key idea: RAII]
The most important idea in C++ has an awkward name: \term{RAII}, ``resource acquisition
is initialisation''. A resource (memory, an open file, a lock, a network connection) is
acquired in an object's constructor and released in its \term{destructor}, which C++ calls
automatically when the object goes out of scope --- even if an exception is thrown.
\code{std::vector}, \code{std::string}, \code{std::unique\_ptr} and \code{std::ifstream}
all work like this. It gives C++ automatic, predictable cleanup without a garbage
collector.
\end{keyidea}

\section{Templates and the Standard Library}

\term[templates]{Templates} let you write code once for many types. The compiler generates a separate,
fully optimised version for each type you use. This is called \term{generic programming},
and it is how the C++ \term{Standard Template Library} (STL) provides containers and
algorithms that work for any element type:
\codefile{C++}{cpp/stl.cpp}{containers, algorithms and a function template}
\outfile{cpp/stl.out}
The STL contains containers (\code{vector}, \code{array}, \code{map},
\code{unordered\_map}, \code{set}, \code{deque}, \code{list}), more than 100 algorithms
(\code{sort}, \code{find}, \code{accumulate}, \code{transform}, \code{count\_if},
\code{minmax\_element}, \ldots), random-number generators (\code{<random>}), time
measurement (\code{<chrono>}), threads (\code{<thread>}), file systems
(\code{<filesystem>}) and complex numbers (\code{<complex>}). Since C++17 many algorithms
can run in parallel by adding one argument, \code{std::execution::par}.

\section{Numerical C++ libraries}

The standard library has no matrices. Engineers use libraries such as:
\begin{itemize}
\item \textbf{Eigen}: a header-only library for vectors, matrices, linear solvers and
  geometry. With \term{operator overloading}, C++ lets the library define what \code{+}
  and \code{*} mean for matrices, so code reads like mathematics:
\begin{lstlisting}[language=C++]
#include <Eigen/Dense>
Eigen::Matrix3d A;
A << 4, -2, 1,
    -2, 4, -2,
     1, -2, 4;
Eigen::Vector3d b(11, -16, 17);
Eigen::Vector3d x = A.ldlt().solve(b);    // x = (1, -2, 3)
\end{lstlisting}
\item \textbf{Armadillo} (MATLAB-like syntax), \textbf{Blaze}, \textbf{xtensor} (NumPy-like);
\item \textbf{PETSc} and \textbf{Trilinos} for large parallel sparse problems;
\item \textbf{Boost}, a large collection of general-purpose libraries, many of which later
  entered the standard;
\item \textbf{OpenMP}, \textbf{MPI}, \textbf{CUDA}, \textbf{SYCL} and \textbf{Kokkos} for
  parallel and GPU computing.
\end{itemize}

\section{Building larger projects}

Real C++ projects consist of many source and header files and depend on external
libraries. Compiling them by hand is impractical. \textbf{CMake} is the de facto standard
build system generator; a minimal \file{CMakeLists.txt} looks like this:
\begin{lstlisting}[language=bash]
cmake_minimum_required(VERSION 3.20)
project(beam LANGUAGES CXX)
set(CMAKE_CXX_STANDARD 20)
add_executable(beam main.cpp sections.cpp)
\end{lstlisting}
\begin{lstlisting}[style=shell]
$ cmake -S . -B build -DCMAKE_BUILD_TYPE=Release
$ cmake --build build
\end{lstlisting}
Package managers such as \textbf{vcpkg} and \textbf{Conan} download and build libraries.

\section{C++ and Python together}

Many \py libraries are C++ underneath, and you can do the same: \textbf{pybind11} and
\textbf{nanobind} make C++ functions and classes importable from \py with a few lines of
glue code. PyTorch, for example, is a C++ library (\emph{LibTorch}) with a \py front end.
This ``\py outside, C++ inside'' pattern is one of the most common architectures in
scientific software.

\section{Strengths and weaknesses}

\begin{center}
\small
\begin{tabularx}{\linewidth}{@{}>{\raggedright\arraybackslash}X>{\raggedright\arraybackslash}X@{}}
\toprule
\textbf{Strengths} & \textbf{Weaknesses}\\
\midrule
Speed of C with high-level abstractions & Very large and complex language\\
Huge ecosystem of engineering and scientific libraries & Long compile times; cryptic template error messages\\
Deterministic resource management (RAII), no garbage-collection pauses & Still possible to make C-style memory errors\\
Runs from microcontrollers to supercomputers and GPUs & No standard package manager or build system\\
Mature tools (debuggers, profilers, IDEs) & Steep learning curve; many outdated tutorials\\
\bottomrule
\end{tabularx}
\end{center}

\section{Further reading}
\begin{itemize}
\item B.~Stroustrup, \emph{A Tour of C++}, 3rd ed., Addison-Wesley, 2022 --- a short
  overview of modern C++ by its creator.
\item B.~Stroustrup, \emph{Programming: Principles and Practice Using C++}, 3rd ed.,
  Addison-Wesley, 2024 --- a first-year university textbook.
\item \url{https://en.cppreference.com} --- the standard reference.
\item The C++ Core Guidelines, \url{https://isocpp.github.io/CppCoreGuidelines/}.
\end{itemize}

\begin{exercises}
\item Rewrite \file{readstats.c} from Chapter~\ref{ch:c} in modern C++. Read the file with
  \code{std::ifstream}, store the values in a \code{std::vector} and add them up with
  \code{std::accumulate}. How much shorter is the program? Is there any \code{free} or \code{delete}?
\item Add an \code{IBeam} class (flange width $b$, total height $h$, flange thickness
  $t_f$, web thickness $t_w$) to \file{sections.cpp}. Compute its second moment of area by
  subtracting two rectangles from the bounding rectangle.
\item Write a function template \code{T clamp\_value(T x, T lo, T hi)} and test it with
  \code{int} and \code{double} arguments. (Then find out that \code{std::clamp} already
  exists.)
\item A list of materials is stored as a \code{std::vector} whose elements are pairs of
  type \code{std::pair<std::string, double>} (material name, density). Use
  \code{std::sort} with a lambda to sort it by density in
  descending order.
\item Use \code{<chrono>} to time a loop that sums $10^8$ numbers. Compare the time with
  and without the \code{-O2} compiler option. What does the optimiser do?
\end{exercises}
